\documentclass[11pt]{article}
\usepackage{amsmath, amssymb}

\begin{document}
	
	\section*{Equivalence of the Definitions of the Second Fundamental Form}
	
	Let $X(u,v)$ be a regular parametrization of a surface $S \subset \mathbb{R}^3$.
	Denote:
	\[
	X_u = \frac{\partial X}{\partial u}, \quad
	X_v = \frac{\partial X}{\partial v}
	\]
	and define the unit normal field
	\[
	N = \frac{X_u \times X_v}{\|X_u \times X_v\|}.
	\]
	
	\subsection*{Definition via Second Derivatives}
	
	The coefficients of the second fundamental form are defined as
	\[
	L = \langle X_{uu}, N \rangle, \quad
	M = \langle X_{uv}, N \rangle, \quad
	N = \langle X_{vv}, N \rangle.
	\]
	
	\subsection*{Definition via the Weingarten Map}
	
	The Weingarten map (shape operator) is defined by
	\[
	dN_p : T_p S \to T_p S.
	\]
	
	Its coefficients are defined via:
	\[
	L = -\langle N_u, X_u \rangle, \quad
	M = -\langle N_u, X_v \rangle = -\langle N_v, X_u \rangle, \quad
	N = -\langle N_v, X_v \rangle.
	\]
	
	\subsection*{Goal}
	
	Show that:
	\[
	\langle X_{ij}, N \rangle = -\langle N_i, X_j \rangle, 
	\quad i,j \in \{u,v\}.
	\]
	
	\subsection*{Proof}
	
	We use the fundamental orthogonality relation:
	\[
	\langle N, X_u \rangle = 0, \quad
	\langle N, X_v \rangle = 0.
	\]
	
	Differentiate $\langle N, X_u \rangle = 0$ with respect to $u$:
	\[
	\frac{\partial}{\partial u} \langle N, X_u \rangle
	= \langle N_u, X_u \rangle + \langle N, X_{uu} \rangle = 0.
	\]
	
	Hence:
	\[
	\langle N, X_{uu} \rangle = - \langle N_u, X_u \rangle,
	\]
	which proves:
	\[
	L = \langle X_{uu}, N \rangle = -\langle N_u, X_u \rangle.
	\]
	
	Similarly, differentiate $\langle N, X_u \rangle = 0$ with respect to $v$:
	\[
	\frac{\partial}{\partial v} \langle N, X_u \rangle
	= \langle N_v, X_u \rangle + \langle N, X_{uv} \rangle = 0,
	\]
	thus:
	\[
	M = \langle X_{uv}, N \rangle = -\langle N_v, X_u \rangle.
	\]
	
	Likewise, differentiating $\langle N, X_v \rangle = 0$ with respect to $u$:
	\[
	\langle N_u, X_v \rangle + \langle N, X_{uv} \rangle = 0,
	\]
	so:
	\[
	M = -\langle N_u, X_v \rangle.
	\]
	
	Finally, differentiate $\langle N, X_v \rangle = 0$ with respect to $v$:
	\[
	\langle N_v, X_v \rangle + \langle N, X_{vv} \rangle = 0,
	\]
	giving:
	\[
	N = \langle X_{vv}, N \rangle = -\langle N_v, X_v \rangle.
	\]
	
	\subsection*{Conclusion}
	
	We have shown:
	\[
	\langle X_{ij}, N \rangle = -\langle N_i, X_j \rangle,
	\]
	thus the two definitions of the second fundamental form coincide.
	
	\hfill $\square$
	
\end{document}
